\chapter{Системийн хөгжүүлэлт}

\section{Хөгжүүлэлтийн орчин, технологийн сонголт}

LexNavigator системийг хөгжүүлэхдээ орчин үеийн веб болон хиймэл оюуны системийн шаардлагад нийцсэн, өргөтгөх боломжтой, өндөр гүйцэтгэлтэй, найдвартай технологиудыг сонгон ашигласан.

Системийн хөгжүүлэлт нь дараах үндсэн зорилгод тулгуурласан:

\begin{itemize}
    \item Хэрэглэгчид ойлгомжтой, хурдан интерфейс бий болгох
    \item Хууль зүйн өгөгдлийг үр ашигтай боловсруулах
    \item AI системтэй уялдсан зөвлөгөө өгөх
    \item Ирээдүйд өргөтгөх боломжтой архитектур бий болгох
\end{itemize}

\noindent Системд ашигласан технологиуд:

\begin{itemize}
    \item \textbf{Frontend:} Next.js, React, Tailwind CSS
    \item \textbf{Backend:} FastAPI, Python
    \item \textbf{Database:} Neon PostgreSQL
    \item \textbf{AI:} Google Gemini 2.5 Flash
    \item \textbf{Security:} JWT, bcrypt
    \item \textbf{DevOps:} Docker, GitHub
\end{itemize}

Эдгээр технологиуд нь системийн гүйцэтгэл, найдвартай байдал болон хөгжүүлэлтийн уян хатан байдлыг хангаж байна.

%----------------------------------------

\section{Системийн модуль, бүрэлдэхүүн хэсгүүд}

LexNavigator систем нь модульчлагдсан бүтэцтэй бөгөөд дараах үндсэн бүрэлдэхүүн хэсгүүдээс бүрдэнэ.

\subsection{Frontend Module}

\begin{itemize}
    \item Хэрэглэгчийн интерфейс
    \item Chat UI
    \item Login / Register хэсэг
    \item AI хариу харах хэсэг
\end{itemize}

\subsection{Backend Module}

Backend нь системийн төв логик хэсэг бөгөөд дараах дэд модулиудаас бүрдэнэ:

\begin{itemize}

    \item \textbf{Auth Module}
    \begin{itemize}
        \item Хэрэглэгчийн бүртгэл, нэвтрэлт
        \item JWT токен үүсгэх, шалгах
    \end{itemize}

    \item \textbf{Chat Module}
    \begin{itemize}
        \item Хэрэглэгчийн асуулт хүлээн авах
        \item AI хариу буцаах
    \end{itemize}

    \item \textbf{Retrieval Module}
    \begin{itemize}
        \item Хууль зүйн өгөгдлөөс холбогдох заалтыг хайх
        \item Context үүсгэх
    \end{itemize}

    \item \textbf{AI Orchestrator}
    \begin{itemize}
        \item Prompt боловсруулах
        \item Gemini API-тай холбогдох
    \end{itemize}

    \item \textbf{Logging Module}
    \begin{itemize}
        \item Хүсэлт, алдаа бүртгэх
        \item Чат түүх хадгалах
    \end{itemize}

\end{itemize}

\subsection{Data Module}

\begin{itemize}
    \item law\_articles -- хууль зүйн заалтууд
    \item categories -- ангилал
    \item user\_chats -- хэрэглэгчийн асуулт, хариу
\end{itemize}

\subsection{Data Processing Module}

\begin{itemize}
    \item PDF хууль унших
    \item Текст цэвэрлэх (Regex)
    \item Өгөгдлийг бүтэцжүүлэх
\end{itemize}

%----------------------------------------

\section{API болон өгөгдөл солилцоо}

LexNavigator систем нь REST API архитектур ашиглан frontend болон backend хооронд өгөгдөл солилцдог.

\subsection{Үндсэн API endpoint-ууд}

\subsubsection{Authentication API}
\begin{itemize}
    \item POST /register -- хэрэглэгч бүртгэх
    \item POST /login -- нэвтрэх
\end{itemize}

\subsubsection{Chat API}
\begin{itemize}
    \item POST /chat -- хэрэглэгчийн асуулт илгээх
\end{itemize}

\subsection{Жишээ хүсэлт}

\begin{verbatim}
{
  "message": "Хулгайн хэрэгт ямар ял оноодог вэ?"
}
\end{verbatim}

\subsection{Жишээ хариу}

\begin{verbatim}
{
  "answer": "Эрүүгийн хуулийн 17.1 дүгээр зүйлд зааснаар..."
}
\end{verbatim}

\subsection{Өгөгдлийн урсгал}

\begin{itemize}
    \item Frontend → Backend (API request)
    \item Backend → Retrieval module
    \item Retrieval → Database
    \item Backend → Gemini API
    \item Backend → Frontend
\end{itemize}

%----------------------------------------

\section{Аюулгүй байдал ба нэвтрэлт}

Системийн аюулгүй байдлыг хангах зорилгоор дараах шийдлүүдийг хэрэгжүүлсэн.

\subsection{Authentication}

\begin{itemize}
    \item JWT (JSON Web Token) ашигласан
    \item Stateless authentication
\end{itemize}

\subsection{Password хамгаалалт}

\begin{itemize}
    \item bcrypt ашиглан нууц үгийг hash хэлбэрээр хадгална
\end{itemize}

\subsection{API хамгаалалт}

\begin{itemize}
    \item Token шалгах middleware
    \item Зөвшөөрөлгүй хүсэлтийг хориглох
\end{itemize}

\subsection{Өгөгдлийн аюулгүй байдал}

\begin{itemize}
    \item SQL injection-ээс хамгаалах (ORM ашиглах)
    \item Input validation (Pydantic)
\end{itemize}

\subsection{Нэмэлт хамгаалалт}

\begin{itemize}
    \item HTTPS ашиглах
    \item Error message-ийг хязгаарлах
\end{itemize}

%----------------------------------------

\section{AI интеграцчилал (Gemini)}

LexNavigator системийн гол онцлог нь хиймэл оюуны интеграцчилал юм.

\subsection{Ажиллах зарчим}

Систем нь хэрэглэгчийн асуултыг шууд AI-д илгээх бус, эхлээд өгөгдлийн сангаас холбогдох хууль зүйн мэдээллийг хайж олсны дараа тэдгээрийг нэгтгэн AI-д дамжуулдаг.

\begin{enumerate}
    \item Хэрэглэгчийн асуултыг backend хүлээн авна
    \item Retrieval module холбогдох хууль зүйн мэдээллийг хайна
    \item Олдсон өгөгдлийг context болгон бэлтгэнэ
    \item Gemini API руу prompt илгээнэ
    \item AI боловсруулсан хариуг буцаана
\end{enumerate}

\subsection{Prompt боловсруулах бүтэц}

\begin{itemize}
    \item System instruction
    \item Хууль зүйн context
    \item Хэрэглэгчийн асуулт
\end{itemize}

\subsection{AI хариу үүсгэх онцлог}

\begin{itemize}
    \item Зөвхөн өгөгдсөн хуульд тулгуурлана
    \item Монгол хэл дээр хариу өгнө
    \item Тайлбарласан, ойлгомжтой хэлбэртэй байна
\end{itemize}

\subsection{AI ашигласны давуу тал}

\begin{itemize}
    \item Хэрэглэгчид хурдан зөвлөгөө өгөх
    \item Хууль ойлгоход хялбар болгох
    \item Мэргэжлийн хэллэгийг энгийн болгох
\end{itemize}

%----------------------------------------

\section{Бүлгийн дүгнэлт}

Энэ бүлэгт LexNavigator системийн хөгжүүлэлт, ашигласан технологи, модуль бүтэц, API харилцаа, аюулгүй байдал болон AI интеграцчиллыг тайлбарлалаа.

Систем нь дараах үндсэн онцлогтой:

\begin{itemize}
    \item Давхаргат, модульчлагдсан архитектур
    \item REST API-д суурилсан өгөгдөл солилцоо
    \item JWT authentication бүхий аюулгүй байдал
    \item AI ашигласан ухаалаг зөвлөгөөний систем
\end{itemize}

Ийнхүү LexNavigator систем нь хууль зүйн мэдээллийг илүү хүртээмжтэй, ойлгомжтой, хурдан хүргэх боломжийг бүрдүүлсэн.